\section{The nested assembly}\label{nested:sec:nested}
\subsection{Input, output, and the scope of the calibration}
An input is a finite real rational $*$-polynomial presentation in typed
centralizer labels and contractions $X_1,\ldots,X_s$. A label has type
projection, unitary, positive contraction, or contraction. The modular
symbols have the prescribed common ratio
\begin{equation}
 h_0X_i=\tfrac12X_ih_0.                                      \label{eq:nested-N0}
\end{equation}
Each monomial contains at most two modular symbols, counting adjoints;
centralizer coefficients can occur before, between, and after them.
Centralizer polynomial equations and degree-one endpoint equations are
included. A normal realization has a nonzero Hilbert space and a faithful
trace-class density $h_0$ of trace one, with all identities on its support.
The ratio in \eqref{eq:nested-N0} is part of the input convention; it is not an operator to be
synchronized. Section~\ref{nested:sec:variants} treats the common ratio $2^{-n}$.
Arbitrary unrelated real ratios are outside this stated lemma: the
calibration used here supplies powers of two.

The proof uses two calibrated levels so that an entire inner activity,
including off-diagonal lifted contractions, is transported by one outer
map. A centralizer corner unitary then identifies equal-density corners.
Products of the resulting modular operators can be translated into actual
centralizer products. The whole-activity refinement is what prevents
independent gauges from changing those products.

\begin{theorem}[Nested compiler]\label{nested:thm:nested}
Every such presentation effectively determines a finite rational Boolean
game. A perfect normal spatial strategy extracts a normal realization
on a source corner of mass $1/81$. Every normal realization has one
simultaneous perfect extension using a $25$-corner scaffold, a qutrit
frame, and a finite row ancilla. The extension preserves finiteness of
every positive moment and of entropy. With types, centralizer conditions,
and modular covariance exact, constructed approximate realizations obey
\begin{equation}
 1-\omega\le C_{\rm nest}\sum_j\|r_j\|_{2,h_0}^2,\qquad
 \|r\|_{2,h_0}^2=\Tr(h_0r^*r),                              \label{eq:nested-N1}
\end{equation}
where $C_{\rm nest}$ depends only on the finite compiled presentation.
\end{theorem}

\subsection{Corners, parameters, and activities}
Begin with nine synchronized corner projections
\[
 f,\ m_{00},m_{01},m_{10},m_{11},\ E_{00},E_{01},E_{10},E_{11}.
\]
Name a synchronized activity ${\Pact_{\rm in}}$ and impose
\begin{equation}
 {\Pact_{\rm in}}=f+\sum_{ab}m_{ab},\qquad {\Pact_{\rm in}}+\sum_{ab}E_{ab}=I.
                                                               \label{eq:nested-N2}
\end{equation}
The inner conditional FMS has flag $f$, activity ${\Pact_{\rm in}}$, and selected corner
$e=m_{00}$. The outer conditional FMS has flag \emph{the whole ${\Pact_{\rm in}}$},
activity $I$, and selected corner $E=E_{00}$. Retain every original
question, both orientations of each event binding, and every inactive
answer. The inner inactive event is $I-{\Pact_{\rm in}}$. The outer inactive event is
zero. In particular its column partition has flag ${\Pact_{\rm in}}$, four $E_{ab}$
answers, and a zero inactive answer. No new inner flag is substituted
for ${\Pact_{\rm in}}$ in that partition.

The frame-bound parameter list is exactly the nine corners, ${\Pact_{\rm in}}$, and
the four positive $ZI,IX$ events, two for each bridge: fourteen
parameters. Use the parameter-binding rows of Section~\ref{sec:projection}.
Set $M_{\rm in}={\Pact_{\rm in}}-f$ and $M_{\rm out}=I-{\Pact_{\rm in}}$. The four unitary axis labels
are equated to $I-2M+2E_{\mathrm{axis},+}$ for their respective bridge.
Let $r_{ab}=x_{\rm in}^az_{\rm in}^be$ and
$R_{ab}=x_{\rm out}^az_{\rm out}^bE$ denote the actual magic corner maps.
These are centralizer expressions in the bound parameters and axes.
The two calibrated transports satisfy
\begin{equation}
 u^*u=f,\ uu^*=e,\ hu=2uh;\qquad
 {\btrans}^*{\btrans}={\Pact_{\rm in}},\ {\btrans}{\btrans}^*=E,\ h{\btrans}=2{\btrans}h.                                 \label{eq:nested-N3}
\end{equation}
The soundness argument below derives the source mass $\Tr(hf)=1/81$ from these relations and the cover.

\subsection{The complete transport list and the whole-activity block}
Encode a projection $P$ by $2P-I$, a positive contraction $C$ by a
unitary with $C=(2I+U+U^*)/4$, a general contraction by the average of
two unitaries, and a unitary directly. Identities in a supported corner
use that corner's identity. The outer transported list consists of:
\begin{enumerate}
\item the two inner projections $f,e$, regarded as projections in ${\Pact_{\rm in}}$;
\item all source label encodings that occur at any endpoint (transporting
all source labels is an allowed deterministic choice);
\item both unitary encodings of every lifted contraction $a_i=e a_i f$;
\item the finitely many derived coefficients specified below when a
coefficient occurs inside a quadratic word.
\end{enumerate}
For every encoding name $k$, introduce global unitaries $U_k,V_k$ with
\[
 (U_k-I)(I-{\Pact_{\rm in}})=0,\qquad (V_k-I)(I-E)=0.
\]
A root-supported encoding additionally fixes $I-f$. Its transported
copy will then fix $E-f'$. Introduce a global unitary ${V^{\rm blk}_k}$ and impose
\begin{equation}
 {V^{\rm blk}_k}={\Pact_{\rm in}} U_k+\sum_{ab}R_{ab}(E V_k)R_{ab}^*.                    \label{eq:nested-N4}
\end{equation}
There is no $I-{\Pact_{\rm in}}-M_{\rm out}$ term since the outer activity is $I$.
In \eqref{eq:nested-N4}, $\Pact_{\rm in}U_k$ acts on the entire five-corner space ${\Pact_{\rm in}}\HH$, including
off-diagonal blocks between $f$ and $e$. Replacing it by a sum of five
unrelated diagonal encodings would lose the required common transport.
The identity $(U_k-I)(I-{\Pact_{\rm in}})=0$ gives $U_k(I-{\Pact_{\rm in}})=I-{\Pact_{\rm in}}$.
Left multiplication by $U_k^*$ gives $U_k^*(I-{\Pact_{\rm in}})=I-{\Pact_{\rm in}}$;
taking adjoints gives $(I-{\Pact_{\rm in}})U_k=I-{\Pact_{\rm in}}$. Thus $[U_k,{\Pact_{\rm in}}]=0$.
The same proof gives $[V_k,E]=0$. Consequently $\Pact_{\rm in}U_k$ is a
unitary on ${\Pact_{\rm in}}\HH$, and each $R_{ab}(EV_k)R_{ab}^*$ is unitary
on its own magic corner. Their supports are orthogonal and cover $I$,
so \eqref{eq:nested-N4} is block diagonal for
${\Pact_{\rm in}}\HH\oplus\bigoplus_{ab}E_{ab}\HH$ and is unitary.

For $f$ use $U_f=I-2{\Pact_{\rm in}}+2f$, and for $e$ use
$U_e=I-2{\Pact_{\rm in}}+2e$; their outer unitary names determine
\[
 f'=\tfrac12 E(I+V_f),\qquad e'=\tfrac12 E(I+V_e).
\]
For each bridge separately refine its retained signs for the six projections
\[
 F_1,\quad F_2,\quad F_3,\quad G_{12}(1),\quad G_{13}(1),\quad G_{23}(1).
\]
For the outer bridge also refine them separately for every $G_{12}({V^{\rm blk}_k})$.
The unchanged opposite questions are retained. There are six refinements at the inner bridge and $6+|\Lambda|$ at
the outer bridge, per retained question and per party.

Label refinement (Lemma~\ref{lem:label}) makes every
retained sign commute with the six frame projections. In the frame
coordinates these generate $\Mat_3\otimes I$: the $F_i$ give the diagonal
matrix units and $F_iG_{ij}(1)F_j=e_{ij}\otimes I/2$ gives the
others. Their commutant is $I_3\otimes B(\HH_{\rm base})$.
Because the corners are frame-bound, the sign and its $(c_0,f)$ block
are frame scalar. Thus ${\btrans}=I_3\otimes\widetilde {\btrans}$ (and likewise $u$).
Writing the outer sign as $I_3\otimes\widetilde{\mathcal X}_0$,
the $(1,2)$ frame block of its commutator with
\[
 G_{12}({V^{\rm blk}_k})=\tfrac12(F_1+F_2+e_{12}\otimes {V^{\rm blk}_k}+
                                      e_{21}\otimes {V^{\rm blk}_k}^*)
\]
is $[\widetilde{\mathcal X}_0,{V^{\rm blk}_k}]/2$. On the outer sign support,
with flag space ${\Pact_{\rm in}}\HH$ and selected magic space $E\HH$, the bridge
lemma gives
\[
 \widetilde{\mathcal X}_0=
 \begin{pmatrix}0&\widetilde {\btrans}^*\\\widetilde {\btrans}&0\end{pmatrix},
 \qquad
 {V^{\rm blk}_k}=\begin{pmatrix}\Pact_{\rm in}U_k|_{{\Pact_{\rm in}}\HH}&0\\0&EV_k|_{E\HH}\end{pmatrix}.
\]
Here $c_0$ means the selected outer magic direction $E$ and $f$
means its flag, which is the \emph{whole} ${\Pact_{\rm in}}$ rather than the inner $f$.
The $(c_0,f)$ block of the commutator is exactly
$\widetilde {\btrans}(\Pact_{\rm in}U_k)-(EV_k)\widetilde {\btrans}$, so
\begin{equation}
                 {\btrans}(\Pact_{\rm in}U_k)=(EV_k){\btrans},                            \label{eq:nested-N5}
\end{equation}
with the same ${\btrans}$ for every $k$, including the two
encodings of a contraction. Here $\Pact_{\rm in}U_k$ may mix all
five inner corners: if $\Pact_{\rm in}U_k=(U_{rs})$ in inner-corner
coordinates, \eqref{eq:nested-N5} says $EV_k={\btrans}(U_{rs}){\btrans}^*$.
Multiplying \eqref{eq:nested-N5} by the unitary adjoints gives transport also for
$U_k^*,V_k^*$. Decoding the same polynomials in these unitaries
therefore transports projections, positive contractions and averages
of two unitaries. Hence $f'={\btrans}f{\btrans}^*$, $e'={\btrans}e{\btrans}^*$,
and, with ${\btrans}_0={\btrans}f$, ${\btrans}_1={\btrans}e$,
\begin{equation}
                  a_i'={\btrans}_1a_i{\btrans}_0^*,                          \label{eq:nested-N6}
\end{equation}
including both encodings of each $a_i$ rather than just its diagonal
compressions. Both sides are actual operators; neither is an
independently gauged copy.

Introduce a centralizer contraction $t$ by two global unitaries and impose
\begin{equation}
 t=f'te,\qquad t^*t=e,\qquad tt^*=f'.                         \label{eq:nested-N7}
\end{equation}
The corner maps $t$ and $a_i$ are synchronized centralizer expressions;
$u,{\btrans},{\btrans}_0,{\btrans}_1$ are extracted transports and are never graph words.
Set $\HH_0=f\HH$ and $h_0=81fhf|_{\HH_0}$. The extracted contractions are
\begin{equation}
                         X_i={\btrans}_0^*ta_i.                      \label{eq:nested-N8}
\end{equation}
The soundness argument below shows that $h_0$ is a faithful trace-one density and that \eqref{eq:nested-N8} satisfies \eqref{eq:nested-N0}.

\subsection{All quadratic translations and arbitrary coefficients}
Put $\ell={\btrans}_0^*t{\btrans}_1^*$. It is a corner unitary from $e'\HH$ to $f\HH$.
Direct multiplication, using only \eqref{eq:nested-N6}--\eqref{eq:nested-N7}, gives
\begin{align}
 X_i^*X_j&=a_i^*a_j,\label{nested:eq:innernew}\\
 {\btrans}_0X_iX_j^*{\btrans}_0^*&=ta_ia_j^*t^*,\label{nested:eq:rangenew}\\
 X_iX_j&=\ell a_i'ta_j.\label{nested:eq:productnew}
\end{align}
For the last identity use $a_i{\btrans}_0^*={\btrans}_1^*a_i'$; the intermediate
projections are $e,f',e'$ and cancel by their stated source/range rules.
The following table is consequently an equivalence in both directions.
\begin{center}\small
\begin{tabular}{@{}ll@{}}
\toprule Root identity & Centralizer identity\\\midrule
$X_i^*X_j={R_{\rm tar}}$ & $a_i^*a_j={R_{\rm tar}}$ at $f$\\
$\sum_iX_iX_i^*={R_{\rm tar}}$ & $\sum_i ta_ia_i^*t^*={R_{\rm tar}}'$ at $f'$\\
$X_iX_j=X_kX_l$ & $a_i'ta_j=a_k'ta_l$\\
$X_i^*X_j=X_kX_l^*$ & $a_i'^*a_j'=ta_ka_l^*t^*$\\
$LX_i=X_jR$ & $L'ta_i=ta_jR$\\\bottomrule
\end{tabular}
\end{center}
Here ${R_{\rm tar}}'={\btrans}_0{R_{\rm tar}}{\btrans}_0^*$ and $L'={\btrans}_0L{\btrans}_0^*$ are obtained from the
transported encodings. The mixed quadratic row uses
${\btrans}_0a_i^*a_j{\btrans}_0^*=a_i'^*a_j'$. Conjugation by ${\btrans}_0$, or multiplication
by $\ell$, is injective on the specified corner. Thus these are
operator equivalences rather than implications inferred from traces.
In particular they give the displayed quadratic identities.

For completeness of the admissible class, here is a finite procedure for
interior coefficients. Give $X$ degree $1$ and $X^*$ degree $-1$.
In a density eigenbasis, a degree-$d$ word has matrix entries only where
the output eigenvalue is $2^{-d}$ times the input eigenvalue. Distinct
degrees are disjoint and orthogonal in $\|\cdot\|_{2,h_0}$; split a
polynomial into degrees $-2,-1,0,1,2$. Take adjoints of negative degrees.
For centralizer coefficients ${K_L},{K_M},{K_R}$:
\begin{enumerate}
\item Degree one: ${K_L}X_i{K_R}$ lifts to ${K_L}'ta_i{K_R}$.
\item Degree zero: ${K_L}X_i{K_M}X_j^*{K_R}$ lifts to
${K_L}'ta_i{K_M}a_j^*t^*{K_R}'$ at $f'$. A term ${K_L}X_i^*{K_M}X_j{K_R}$ first becomes
${K_L}a_i^*t^*{K_M}'ta_j{K_R}$ at $f$; transport that whole supported centralizer
expression to $f'$ through its own derived encoding.
\item Degree two: put $H_{K_L}=t^*{K_L}'t$ at $e$ and transport it to
$H_{K_L}'={\btrans}_1H_{K_L}{\btrans}_1^*$ at $e'$. Then
\[
             {K_L}X_i{K_M}X_j{K_R}=\ell H_{K_L}'a_i'{K_M}'ta_j{K_R}.
\]
Indeed $\ell H_{K_L}'={K_L}\ell$, by multiplication and $tt^*=f'$.
\end{enumerate}
A derived coefficient with a known rational norm bound $K$ is encoded
after division by $\max(1,K)$, and its defining equation is added as a
centralizer polynomial. Such a bound follows by multiplying and summing
the declared contraction bounds. Only the two displayed kinds of derived
coefficients are needed for degree two. Define the base labels and $t$
first, the original transports next, and then the derived coefficients
and their transports. Definitions involve no newly transported derived
coefficient. Adding their
separate refinements constrains the original
${\btrans}$, and never extracts a fresh gauge. This is a finite, non-circular
assembly for the whole stated class.

\subsection{Soundness and the Boolean game}
Start with a normal strategy, purify without querying the environment,
and reduce to the original marginal supports by facts (N1)--(N8) of Section~\ref{sec:compiler}. Perfect matching precedes every use of projectivity. Boolean
row tests give the rational identities on these supports. Graphs and
triangles then give actual common unitary words. Parameter bindings
identify the actual FMS events; they do not replace them by free labels.
Apply the two conditional bridge theorems to those bound events. The
frame refinements apply first; the block refinements then give \eqref{eq:nested-N5}.
Decode labels, impose \eqref{eq:nested-N7}, and apply \eqref{eq:nested-N8} and the translation procedure.
This order uses neither the intended application nor its spectral
classification in extracting the scaffold.

All cancellations of a coefficient operator use its dense range on the
original local support, as proved in (N1)--(N8). On a mixed-state cut
$A\mid BE$ the coefficient map need not be injective. All graph calculations are
bounded block identities. The source density is faithful since $f$ commutes with $h$ and
$h$ is faithful on the reduced support.

Write $h=h_A$ on $P_A\HH_A$. The synchronized
corners are projections commuting with $h$; their sum identities
make them orthogonal. Since $hu=2uh$ and $h{\btrans}=2{\btrans}h$, cyclicity of the
trace of a trace-class operator times bounded operators gives
\[
 \begin{aligned}\Tr(he)&=\Tr(huu^*)=\Tr(u^*hu)=2\Tr(hf),\\
 \Tr(hE)&=\Tr(h{\btrans}{\btrans}^*)=\Tr({\btrans}^*h{\btrans})=2\Tr(h{\Pact_{\rm in}}).\end{aligned}
\]
The centralizer magic-corner maps give equal masses to the four
inner selected copies and, separately, the four outer selected
copies. Thus $\Tr(h{\Pact_{\rm in}})=9\Tr(hf)$ and
$1=\Tr(h{\Pact_{\rm in}})+4\Tr(hE)=9\Tr(h{\Pact_{\rm in}})$, proving $\Tr(hf)=1/81$.
The bridge lemma holds also for a zero activity, so no nonvanishing sector
is assumed here. It follows that $\HH_0=fP_A\HH_A\ne0$ and
$h_0=81fhf|_{\HH_0}$ is faithful, trace class and trace one.

A supported unitary commutes with its support by the calculation
above. On $\HH_0$ the four decoding types are therefore respectively
$U$, $(I+U)/2$ with $U^2=I$, $(2I+U+U^*)/4$, and
$(U_++U_-)/2$. They are respectively unitary, projection, positive
contraction, and contraction, and each commutes with $h_0$.
The supported $a_i=e a_i f$ is a contraction; \eqref{eq:nested-N7} makes $t$ a
partial isometry from $e$ onto $f'$, and ${\btrans}_0$ is one from $f$ onto
$f'$. Hence $X_i={\btrans}_0^*ta_i$ maps $\HH_0$ into itself with norm at
most one. Since $h{\btrans}_0^*={\btrans}_0^*h/2$ and all other factors centralize $h$,
\[
 h_0X_i=81fh{\btrans}_0^*ta_i
 =\tfrac{81}{2}{\btrans}_0^*hta_i
 =\tfrac{81}{2}{\btrans}_0^*ta_ifhf=\tfrac12X_ih_0.
\]
Transport of decoded labels and their adjoints, followed by the
bounded corner multiplications in Section~\ref{nested:sec:nested}, now gives every
imposed application identity on this actual $\HH_0$ in both
directions. The degree splitting uses the density eigenbasis on
this support, never an inverse of $h$.

Bob's original marginal support is treated in the same way: matched
coarse effects are its projective mirrors, the bridge and refinement
lemmas apply on that support, and the same cyclicity gives source
mass $1/81$. Mirroring products reverses their order, as in (N4a);
no extra mirror equation is imposed. Alice's extracted realization
already proves soundness; the symmetric support argument justifies
all Bob-side bindings and refinements, including on a purified
mixed-state cut with its original Bob marginal.

Expand the final finite centralizer polynomials in their real rational
unitary words, cancel adjacent inverse tokens, and close under proper
prefixes.
Add all graphs, placement triangles, product triangles,
inverse triangles, parameter bindings, and the off-diagonal polynomial
rows. Apply the Boolean-row construction already proved above. Retain
all FMS and refinement questions, pad the two question sets, and sample
every pair uniformly with positive rational probability $Q^{-2}$.
Accept every unspecified answer pair. Perfection gives zero probability
for each individual rejecting rule, proving soundness.

\subsection{One simultaneous completeness model}
Let $\mu=\operatorname{diag}(1,2,2,2,2)/9$. On
$\mathbb C^5_{\rm out}\otimes\mathbb C^5_{\rm in}\otimes\HH_0$ take
\begin{equation}
                     h_{\rm sc}=\mu\otimes\mu\otimes h_0.     \label{eq:nested-N9}
\end{equation}
Outer index $0$ is the inner activity; inner index $0$ is $f$ and
inner index $1$ is $e$. The selected outer index is $1$. Thus
\begin{align*}
 f&=|0,0\rangle\langle0,0|\otimes I,&
 e&=|0,1\rangle\langle0,1|\otimes I,\\
 f'&=|1,0\rangle\langle1,0|\otimes I,&
 e'&=|1,1\rangle\langle1,1|\otimes I,\\
 {\btrans}&=|1\rangle\langle0|\otimes I_5\otimes I,&
 t&=|1,0\rangle\langle0,1|\otimes I,\\
 a_i&=|0,1\rangle\langle0,0|\otimes X_i,&
 a_i'&=|1,1\rangle\langle1,0|\otimes X_i.
\end{align*}
The density weights at $e$ and $f'$ are both $2h_0/81$, so $t$
centralizes $h_{\rm sc}$. Commutation of $a_i$ with $h_{\rm sc}$ is
exactly $2h_0X_i=X_ih_0$. All the copies use the identity outer
identifications. The inner and outer FMS effects are the real canonical
ones on their respective five-dimensional factors, with the inner
copy inactive off outer index $0$. Their compressed sign vectors are
$|f,f\rangle+(-1)^{xy}|c_x,r_y\rangle$, up to the common positive
amplitude and spectator factors. The FMS sign completeness used here
is the explicit calculation at the end of the proof of
Lemma~\ref{lem:four}.

Every positive density eigenspace is finite-dimensional. For a
centralizer contraction $T$, its polar partial isometry extends on
each such eigenspace to a unitary $V$. Then
\begin{equation}
 U_\pm=V\bigl(|T|\pm i\sqrt{I-|T|^2}\bigr),\qquad
                         T=(U_++U_-)/2.                     \label{eq:nested-N10}
\end{equation}
Direct sums give bounded centralizer unitaries. On a prescribed
support such as ${\Pact_{\rm in}}$, do this inside that support and extend by identity
outside. For a root positive contraction use
$U=2C-I+2i\sqrt{C(I-C)}$, and for a projection use $2P-I$.
Encode source choices once and transport precisely those choices into
$E$; \eqref{eq:nested-N4} now holds simultaneously for every label, even a lifted
$a_i$ with off-diagonal inner blocks. Encode the global $t$ by \eqref{eq:nested-N10}.

Append $I_3/3$ and all graph, triangle, and parameter-binding
projections. For every refinement use actual joint products: its
block unitary commutes with its own FMS signs. Distinct labels require
no mutual commutation. Finally append $I_{d_{\max}}/d_{\max}$ and the
Boolean-row effects. Alice uses these operators and Bob their
transposes in a product density eigenbasis. Facts (N1)--(N8) prove every
matching, calibration, orthogonality, and row-zero test with this one
state and assignment.

Both nested bridges have magic density twice the
flag density. For the inner one the flag spectator density is $h_0$
(up to its scalar outer weight); for the outer one it is
$\mu\otimes h_0$. Denote either normalized spectator by $\rho$,
and the total flag mass by $\lambda$. The canonical magic directions
satisfy $\langle c_x,r_y\rangle=\epsilon_{xy}/\sqrt2$, where
$\epsilon_{xy}=(-1)^{xy}$. Compressing $\Omega_h$ to the two sign
supports gives, omitting only the unchanged frame and row factors,
\[
 \Omega_{xy}=\sqrt\lambda\bigl(|f,f\rangle+
            \epsilon_{xy}|c_x,r_y\rangle\bigr)\otimes\Omega_\rho.
\]
Alice's sign observable exchanges $f,c_x$ and Bob's exchanges $f,r_y$.
Applying the swaps gives
$(\mathcal X_x\otimes I)\Omega_{xy}
=\epsilon_{xy}(I\otimes\mathcal Y_y)\Omega_{xy}$.
Their spectral sign answers therefore never violate $ab=\epsilon_{xy}$.
The other column/row projections give the retained answers and the
calibration rules exactly. The inner measurements have inactive
effect $I-{\Pact_{\rm in}}$ off outer index $0$; that sector is matched on the two
sides. The outer bridge is active everywhere.

Under $E\HH\cong\mathbb C^5_{\rm in}\otimes\HH_0$, let
$\widehat U_k$ be the entire chosen unitary on ${\Pact_{\rm in}}\HH$.
Identity outer identifications make
\[
                    {V^{\rm blk}_k}=I_{\rm out}\otimes\widehat U_k.
\]
Every outer FMS effect, including its noncentralizer sign effects,
acts on the outer factor alone and hence commutes with ${V^{\rm blk}_k}$,
even when $\widehat U_k$ has off-diagonal inner blocks. The inner
bridge has only the six frame refinements, which commute because
the frame is a separate scalar factor. Thus every labelling
refinement is the actual positive joint product of commuting
projections, sums to the original coarse measurement, and has the
required graph marginal. Transposition in the product density eigenbasis preserves
these commutators and products on Bob's side. All centralizer
matching and orthogonality tests follow from \eqref{eq:nested-N8}, and the appended
row factor gives all Boolean row tests simultaneously with the same
state. This proves the claimed nested completeness assignment.

The complete marginal is a tensor product of $h_0$ with a fixed
finite-dimensional density $\eta$. Thus
$\Tr[(\eta\otimes h_0)^\alpha]=\Tr(\eta^\alpha)\Tr(h_0^\alpha)$
for every $\alpha>0$, and its entropy is $S(\eta)+S(h_0)$.
Conversely compression to the commuting source corner of mass $1/81$
cannot create a finite positive moment from an infinite full moment.
These observations give the stated resource preservation.

\subsection{The loss constant, including the normalization factors}
Keep exact structural hypotheses and choose all derived coefficient
encodings from their actual expressions, so their defining rows stay
exact. A homogeneous source residual $r$ has a lifted centralizer
residual $p$. For a degree-one or degree-two lowering residual the
lift has input corner $f$, so
\[
 \|p\|_{2,h_{\rm sc}}^2=\frac1{81}\|r\|_{2,h_0}^2.
\]
A degree-zero residual transported to $f'$ instead has factor $2/81$.
For negative degrees take the adjoint; if $h_0r=2^{-d}rh_0$,
cyclicity gives $\|r^*\|_{2,h_0}^2=2^{-d}\|r\|_{2,h_0}^2$.
Therefore each chosen translation has factor $\beta_j/81$, with
$\beta_j\le4$. Splitting degrees preserves the sum of squared norms
because their density matrix entries have disjoint ratios.

For the qutrit row $\Delta_j$ the two off-diagonal blocks are $p_j,p_j^*$.
Since $p_j$ is in the centralizer, trace cyclicity gives
\[
 \Tr(h_{\rm old}\Delta_j^2)=\frac{2\beta_j}{243}\|r_j\|_{2,h_0}^2.
\]
Let $L_j$ clear the rational row denominators and let $d_j$ be its
padded power-of-four column dimension. Equation~\eqref{eq:rowloss}
then gives, for one oriented row-zero test,
\begin{equation}
 Q^{-2}p(\Pi^{\rm row}_0,\Pi^{\rm ref}_{d_j})=
 \frac{8L_j^2\beta_j}{243d_j^3Q^2}\|r_j\|_{2,h_0}^2.
                                                               \label{eq:nested-N11}
\end{equation}
Sum these probabilities; overlapping rules can only reduce total loss.
If both orientations are retained, include both in the sum (equivalently
multiply the maximum constant by two). One may therefore take the
maximum displayed coefficient, with the explicit orientation
multiplicity, as $C_{\rm nest}$. It is independent of the Hilbert space,
density spectrum, witness period, and computation time. If just one
forbidden projection has root probability $p$, its residual squared norm
is $p$, and \eqref{eq:nested-N1} reads loss at most $C_{\rm nest}p$. This proves the
completeness estimate.

\section{The bridge-count and branching variants}\label{nested:sec:variants}
\begin{corollary}[Common dyadic ratio and arbitrary finite branching]
For every fixed $n\ge1$, replace the inner bridge by an $n$-bridge chain
and the outer bridge by an $n$-bridge chain with first flag the entire
inner activity. The same degree-two compiler works with common ratio
$2^{-n}$. If $d_n=1+4n$ and
$Z_n=1+4\sum_{g=1}^n2^g$, the scaffold has $d_n^2$ scalar corners and
source mass $Z_n^{-2}$. Any finite number $b$ of creation labels may
be used. For a full $b$-ary isometry presentation a normal realization
requires $b2^{-n}<1$ and contains the weighted Fock factor with weights
$(1-b2^{-n})2^{-n|w|}$.
\end{corollary}
\begin{proof}
The explicit chain corner list is $e_0,m_{g,ab}$, $1\le g\le n$,
with $e_g=m_{g,00}$, activity $e_{g-1}+\sum_{ab}m_{g,ab}$, and cover
$e_0+\sum_{g,ab}m_{g,ab}=I$. The bridge transports ${\btrans}_g$ have ratio
$2$; their product has ratio $2^n$. Use the block equation \eqref{eq:nested-N4} at
each link, including its inactive identity term when the link activity
is not $I$. Transport the same encoding list through every link.
This gives the actual common product transport, by induction on $g$.
The inner chain density is
\[
 Z_n^{-1}\operatorname{diag}(1;2^{\oplus4};4^{\oplus4};\ldots;
 (2^n)^{\oplus4})\otimes h_0.
\]
Take two such factors. The inner last selected corner and the outer
last selected copy of the source both have density $2^nh_0/Z_n^2$,
so the same $t$ construction applies. Replace ${\btrans}_0,{\btrans}_1$ by restrictions
of the common outer chain product. Every multiplication in the proof
above remains valid, and the common quadratic factor still exists.
The finite mass factors in \eqref{eq:nested-N11} change accordingly and remain
presentation-dependent constants.

The number $b$ affects only the finite list of lifted $a_i$ and its
isometry, orthogonality, and range rows. No bridge has to possess $b$
magic corners: its magic square still has four. If $A_i^*A_j=\delta_{ij}I$
and $hA_i=2^{-n}A_ih$, cyclicity gives total range mass $b2^{-n}$.
Pull a positive density eigenvector backwards through nonzero $A_i^*$;
its eigenvalue increases by $2^n$ and so the process must terminate.
The defect and all its word descendants exhaust the support. The word
sectors are orthogonal, and their densities are $2^{-n|w|}$ times
the root density. Normalization forces $b2^{-n}<1$ and gives the
stated factor. Its $\alpha$-moment converges precisely when
$b2^{-n\alpha}<1$. This proves both the $n\ne4$ and $b\ne4$ variants.
\end{proof}

\paragraph{Exactly $n$ bridges for the bare-tree presentations of Section~\ref{sec:renyi}.}
The bare-tree and degree-one endpoint presentations need only one
chain, hence exactly $n$ bridges, rather than the two chains needed
for the quadratic extension. Here is its complete specialization.
Use the preceding $1+4n$ corner list and cover, bind those corners and
the $2n$ positive axis events to the frame, and set
$v={\btrans}_n\cdots {\btrans}_1$. For each transported encoding use the
copies $u_{g,\ell}$ fixing $I-e_g$ and the block equation
\eqref{eq:blockcopies}. Refine that link for the six frame projections and every
$G_{12}(V_{g,\ell})$. This gives the same calibrated transport of every
label through $v$, with no independent link gauges. Introduce
$b_i=e_nb_ie_0$ and extract $A_i=v^*b_i$. The input rows are exactly
\[
 b_i^*b_j=\delta_{ij}e_0,\qquad
 \sum_i b_ib_i^*=e_n-P_n,\qquad L_nb_i=b_jR_0,
\]
plus the root centralizer polynomial rows. These are respectively
the isometry/orthogonality, range, and endpoint translations; multiply
by $v,v^*$ to obtain both implications. The cover fixes source mass
$Z_n^{-1}$. The one-chain density displayed in the proof, the same
finite-block encodings, and actual joint refinements give simultaneous
completeness. The one-oriented loss constant is
\[
 \max_j\frac{8L_j^2\beta_j}{3Z_nd_j^3Q^2},\qquad
 \beta_j=1\ \hbox{for source and endpoint rows},\quad
 \beta_j=2^n\ \hbox{for range rows}.
\]
For the two-chain quadratic version replace $Z_n$ by $Z_n^2$ and use
$\beta_j\le2^{2n}$. Any finite $b$ is handled just by the displayed
finite $i,j$ lists. Thus the one-chain $n$-bridge compiler used in Section~\ref{sec:renyi} and the nested degree-two extension are both specified.

